\section{ماژول tomoe\_compressor (ToMoE Compressor)}
\label{sec:module_tomoe_compressor}

\subsection{معرفی و انگیزه}
این مکانیزم بر مبنای مبانی نظری پیشرفته در ادبیات توجه و فشرده‌سازی \cite{eckart1936approximation} توسعه یافته است.

تبدیل لایه‌های متراکم خطی به ساختار آمیخته‌ای از خبره‌ها (\lr{Mixture of Experts - MoE}) یکی از نوآورانه‌ترین راهکارها برای حفظ ظرفیت مدل در عین افزایش کارایی محاسباتی و کاهش بار محاسباتی فعال در هر گام استنباطی است. مدل‌های سنتی ترنسفورمر حجم عظیمی از محاسبات را در لایه‌های پیش‌خوران (\lr{FFN}) و نگاشت‌های پروجکشن توجه صرف می‌کنند که در آن‌ها تمامی نورون‌ها برای هر توکن ورودی فعال می‌شوند.

موتور فشرده‌سازی \lr{ToMoE Compressor (Top-k Mixture of Experts by Weight Clustering)} با هدف بازآرایی ساختاری شبکه‌های عصبی چگال به معماری خلوت و آگاه از مسیر (\lr{Sparse Routed Architecture}) توسعه یافته است. در این الگوریتم، به جای آموزش یک مدل \lr{MoE} از نقطه صفر، وزن‌های لایه‌های خطی و \lr{Conv1D} از طریق خوشه‌بندی نورون‌های خروجی (\lr{Weight Clustering}) به چند زیرمجموعه متمایز از خبره‌ها افراز شده و یک شبکه مسیریاب پویا (\lr{Router}) وظیفه تخصیص توکن‌ها به برترین خبره‌ها (\lr{Top-k Routing}) را بر عهده می‌گیرد.

\subsection{مبانی نظری و فرمول‌بندی ریاضی}
فرض کنید $\mathbf{W} \in \mathbb{R}^{d_{\text{out}} \times d_{\text{in}}}$ ماتریس وزن یک لایه خطی متراکم باشد که می‌توان آن را به صورت مجموعه‌ای از بردارهای ردیفی (نورون‌های خروجی) نمایش داد:
\begin{equation}
\mathbf{W} = \begin{bmatrix} \mathbf{w}_1^\top \\ \mathbf{w}_2^\top \\ \vdots \\ \mathbf{w}_{d_{\text{out}}}^\top \end{bmatrix}, \quad \mathbf{w}_i \in \mathbb{R}^{d_{\text{in}}}
\end{equation}

الگوریتم \lr{ToMoE} با اجرای خوشه‌بندی \lr{K-Means} با دقت محاسباتی \lr{FP32} (یا \lr{MiniBatchKMeans} برای لایه‌های بسیار بزرگ با $d_{\text{out}} > 2048$) نورون‌ها را به $E$ خوشه متمایز $\mathcal{C}_1, \mathcal{C}_2, \dots, \mathcal{C}_E$ با ابعاد ترجیحاً برابر افراز می‌نماید:
\begin{equation}
\arg\min_{\{\mathcal{C}_e\}_{e=1}^E} \sum_{e=1}^E \sum_{\mathbf{w}_i \in \mathcal{C}_e} \|\mathbf{w}_i - \boldsymbol{\mu}_e\|_2^2
\end{equation}
که در آن $\boldsymbol{\mu}_e$ مرکز خوشه $e$-ام است. در صورت تولید خوشه‌های خالی یا نابرابری شدید، مکانیزم پشتیبان افراز مساوی و معین (\lr{Deterministic Equal Division}) فعال می‌شود.

هر خوشه متناظر با یک خبره $\mathbf{W}^{(e)} \in \mathbb{R}^{d_e \times d_{\text{in}}}$ است. برای هر بردار ورودی $\mathbf{x} \in \mathbb{R}^{d_{\text{in}}}$، شبکه مسیریاب با استفاده از نگاشت خطی $\mathbf{W}_{\text{gate}} \in \mathbb{R}^{E \times d_{\text{in}}}$ امتیازهای هدایت را در فضای \lr{FP32} محاسبه می‌کند:
\begin{equation}
\mathbf{h} = \mathbf{x} \mathbf{W}_{\text{gate}}^\top + \mathbf{b}_{\text{gate}}
\end{equation}
سپس، عملگر انتخاب برترین شاخص‌ها (\lr{Top-k}) وزن‌های ترکیب را تعیین می‌نماید:
\begin{equation}
\mathcal{T} = \operatorname{Top-k}(\mathbf{h}, k), \quad g_e = \begin{cases} \frac{\exp(h_e)}{\sum_{j \in \mathcal{T}} \exp(h_j)} & \text{if } e \in \mathcal{T} \\ 0 & \text{otherwise} \end{cases}
\end{equation}

خروجی موثر نهایی از پراکنده‌سازی و بازنشانی خروجی خبره‌های فعال بر روی بردار کامل بعد خروجی بازسازی می‌شود:
\begin{equation}
\mathbf{y} = \sum_{e \in \mathcal{T}} g_e \cdot \operatorname{Scatter}\left(\mathbf{x} (\mathbf{W}^{(e)})^\top + \mathbf{b}^{(e)}, \mathcal{I}_e\right)
\end{equation}
که در آن $\mathcal{I}_e$ اندیس نورون‌های متعلق به خوشه $e$ در ماتریس خروجی اصلی است.

\subsection{معماری و پیاده‌سازی}
کلاس‌های اصلی این موتور شامل \lr{ToMoELinear} و \lr{ToMoECompressor} بوده و در رجیستری فشرده‌سازی سیستم تحت نام‌های \lr{tomoe} و \lr{tomoe\_compressor} به ثبت رسیده‌اند. شاخص‌ترین قابلیت‌های معماری و فنی این پیاده‌سازی عبارتند از:
\begin{itemize}
    \item \textbf{مسیر پیش‌رو کاملاً برداری‌شده (\lr{Fully Vectorized Forward Pass}):} پیاده‌سازی بدون هیچ‌گونه حلقه تکرار پایتونی روی توکن‌ها انجام شده و با استفاده از دستورات \lr{torch.gather} و عملگرهای \lr{scatter\_add\_} برای خبره‌های هم‌اندازه به بیشترین بازده اجرایی در کارت‌های گرافیکی دست می‌یابد.
    \item \textbf{پایداری محاسباتی در دقت‌های مختلط:} محاسبات امتیازدهی و توابع بیشینه هموار (\lr{Softmax}) در شبکه مسیریاب در فضای \lr{FP32} انجام گرفته و همگامی کامل با قالب‌های داده \lr{FP16} و \lr{BF16} حفظ می‌شود.
    \item \textbf{انطباق با نگاشت‌های دوگانه و مدرن:} پشتیبانی جامع از وزن‌های ترانهاده \lr{Conv1D} در مدل‌های خانواده \lr{GPT-2} و پروجکشن‌های ذوب‌شده و \lr{GQA} در مدل‌های نظیر \lr{Qwen/LLaMA}.
    \item \textbf{محافظت از لایه‌های حیاتی:} لایه‌های حساس شامل سر دسته‌بندی (\lr{classifier})، سر مدل زبانی (\lr{lm\_head})، ادغام‌کننده (\lr{pooler}) و لایه‌های تعبیه (\lr{embedding}) به منظور پیشگیری از فروپاشی شناختی مدل از فرآیند تبدیل به \lr{MoE} مستثنی می‌شوند.
    \item \textbf{پشتیبانی اختیاری از تطبیق‌دهنده‌های کم‌رتبه (\lr{LoRA}):} امکان تعریف پارامترهای آموزش‌پذیر \lr{LoRA} بر روی خبره‌ها و شبکه مسیریاب برای بازیابی سریع و دقیق عملکرد مدل پس از تبدیل ساختاری.
\end{itemize}

\subsection{نتایج تجربی}
ارزیابی بنچ‌مارک سخت‌افزاری و آزمون دقت فشرده‌ساز \lr{tomoe\_compressor} بر روی مدل‌های زبانی گوناگون بر پایه نتایج مستند ارزیابی‌های جامع در جدول \ref{tab:benchmark_tomoe_compressor} گزارش شده است:

\begin{table}[H]
\centering
\small
\caption{نتایج تجربی و اعتبارسنجی بنچ‌مارک فشرده‌ساز \lr{tomoe\_compressor} روی مدل‌های پایه}
\label{tab:benchmark_tomoe_compressor}
\begin{tabular}{lccccc}
\toprule
\textbf{مدل پایه} \lr{(Base Model)} & \textbf{دقت} \lr{(Accuracy)} & \textbf{شاخص} \lr{(F1-Score)} & \textbf{حافظه} \lr{(Peak VRAM)} & \textbf{توان} \lr{(Throughput)} & \textbf{وضعیت} \\
\midrule
\lr{DistilBERT} & \lr{69.50\%} & \lr{0.7569} & \lr{427.00 MB} & \lr{9,828.24 tok/s} & \lr{PASS} \\
\lr{RoBERTa} & \lr{50.92\%} & \lr{0.6748} & \lr{571.84 MB} & \lr{5,169.98 tok/s} & \lr{PASS\_LOW} \\
\lr{GPT-2} & \lr{54.24\%} & \lr{0.6780} & \lr{874.64 MB} & \lr{6,994.41 tok/s} & \lr{PASS\_LOW} \\
\lr{Qwen} & \lr{---} & \lr{---} & \lr{---} & \lr{---} & \lr{NOT\_RUN} \\
\bottomrule
\end{tabular}
\end{table}

بر اساس نتایج تجربی، ماژول \lr{tomoe\_compressor} در مدل \lr{DistilBERT} موفق به دستیابی به دقت پایدار \lr{69.50\%} و شاخص F1 معادل \lr{0.7569} با وضعیت معتبر \lr{PASS} شده است. در مدل‌های بزرگ‌تر \lr{RoBERTa} و \lr{GPT-2} به دلیل حساسیت شدید بردارهای بازنمایی پیش‌خوران به تغییرات ناگهانی خلوت‌سازی بدون آموزش طولانی‌مدت مسیریاب، دقت مدل به ترتیب \lr{50.92\%} و \lr{54.24\%} ثبت شده که با برچسب \lr{PASS\_LOW} علامت‌گذاری شده است. مدل \lr{Qwen} در فاز بنچ‌مارک متمرکز ارزیابی نشده و وضعیت آن به عنوان \lr{NOT\_RUN} ثبت شده است.

\subsection{محدودیت‌ها و تحلیل شکست}
تحلیل دقیق رفتار الگوریتم و شکست‌های مشاهده‌شده حاکی از نکات زیر است:
\begin{itemize}
    \item \textbf{افت دقت صفر-گذر (\lr{Zero-shot Degradation}) بدون ریزتنظیم مسیریاب:} از آنجا که وزن‌های مسیریاب به صورت تصادفی با توزیع نرمال کوچک مقداردهی اولیه می‌شوند، عدم وجود گام‌های کافی برای آموزش مسیریاب سبب می‌شود که توکن‌ها به خبره‌های نامرتبط هدایت گردند و دقت اولیه مدل افت نماید (مانند افت به \lr{50.92\%} در \lr{RoBERTa}).
    \item \textbf{عدم تقارن خوشه‌ها در وزن‌های نامنظم:} با وجود استفاده از مکانیزم بازتخصیص متوازن، توزیع وزن‌های پیش‌خوران لزوماً ساختار خوشه‌ای طبیعی کروی ندارند که این امر خطای تقریب بازنمایی نورون‌ها را افزایش می‌دهد.
    \item \textbf{بار اضافه حافظه ناشی از چند خبره همزمان:} در صورتی که تمامی خبره‌ها در حافظه گرافیکی نگهداری شوند و اندازه بافر اندیس‌ها بزرگ باشد، مقدار مصرف حافظه لحظه‌ای به دلیل پراکنده‌سازی تانسورها اندکی افزایش می‌یابد.
\end{itemize}
